% Generated draft: requires mathematical and source review.
\documentclass[11pt]{article}
\usepackage[margin=25mm]{geometry}
\usepackage{fontspec}
\setmainfont{Latin Modern Roman}
\usepackage{amsmath,amssymb,mathtools,unicode-math}
\setmathfont{Latin Modern Math}
\usepackage{enumitem,needspace,xurl,hyperref}
\setlength{\parindent}{0pt}
\setlength{\parskip}{5pt}
\setlength{\emergencystretch}{3em}
\newcommand{\N}{\mathbb{N}}
\newcommand{\Z}{\mathbb{Z}}
\newcommand{\Q}{\mathbb{Q}}
\newcommand{\R}{\mathbb{R}}
\newcommand{\C}{\mathbb{C}}
\newcommand{\F}{\mathbb{F}}
\newcommand{\A}{\mathbb{A}}
\newcommand{\PP}{\mathbb{P}}
\newcommand{\E}{\mathbb{E}}
\newcommand{\eps}{\varepsilon}
\newcommand{\dd}{\,\mathrm{d}}
\begin{document}
\section*{UnsolvedMath ID: 30004839}
\textbf{Source problem code:} catalogue entry
\subsection{Definitions and mathematical statement}
For terminology, let \(A=\operatorname{Pol}(G_q)\) be the Hopf \(*\)-algebra of representative functions of a compact quantum group \(G_q\), with counit \(\varepsilon:A\to\C\), and put \(K=\ker\varepsilon\). A generating functional is a linear map \(L:A\to\C\) satisfying
\[
L(1)=0,\qquad L(a^*)=\overline{L(a)}\quad(a\in A),\qquad L(a^*a)\geq0\quad(a\in K).
\]
A generating functional is called Gaussian if
\[
L(abc)=0\qquad(a,b,c\in K).
\]
These standard definitions provide terminology; they do not determine the particular decomposition method intended by the source.

\textbf{Source conjecture.} Does the L\'evy--Khintchine decomposition method described in the cited report for the quantum group denoted there by \(U_q(n)\) extend to the \(q\)-deformations \(G_q\) of other simple compact Lie groups \(G\)?

The intended assertion concerns extending that particular method across other simple compact Lie groups. The supplied record does not specify the admissible values of \(q\), the precise construction of \(G_q\), any restrictions on the global form of \(G\), the class of functionals covered, or the exact decomposition and its existence or uniqueness requirements. Consequently, a fully quantified decomposition conjecture cannot be recovered from this record alone. In particular, no particular decomposition formula or universal range of \(q\) is asserted here.

The original extraction is spliced with auxiliary material. It includes an assertion involving the undefined abbreviation ``GFs'', a formula resembling
\[
\psi(a)=\langle v,[\pi(a)-\varepsilon(a)I]v\rangle,
\]
and a statement that each \(\eta_j\) is determined by a value at \(u_{jj}\), followed by a limit involving \(\pi(a)-\varepsilon(a)I\), \((I-p\pi(u_{jj}))^{-1}\), and \(p\to1\). The extraction alternates between \(\eta(u_{jj})\in H_j\) and \(\eta_j(u_{jj})\), and supplies no definitions or hypotheses for these expressions. These fragments do not specify additional claims of the present conjecture; their intended scope and relation to the method require the original context.
\subsection{Short English statement}
The question asks whether the report's L\'evy--Khintchine decomposition method for \(U_q(n)\) extends to quantum deformations of other simple compact Lie groups. The parameter range, functional class, and exact decomposition are unspecified. The supplied dated assessment labels the problem partially solved; this draft does not independently verify that assessment.
\subsection{Sources}
\begin{itemize}
\item Local UnsolvedMath registry: \texttt{data/unsolvedmath/problems.json}, permanent ID \(30004839\), source problem code \texttt{catalogue entry}, title ``Gaussian Functionals on q-Deformed Compact Lie Groups''. The supplied record names no proposer and records proposed year \(2021\).
\item Supplied primary citation: \emph{Quantum Groups -- Algebra, Analysis and Category Theory} (2022), Oberwolfach Reports, pp.\ 2397--2458, DOI \texttt{10.4171/owr/2021/44}. Source URL and literature-source label ``Oberwolfach Report 44'': \url{https://doi.org/10.4171/owr/2021/44}. The supplied background describes the reports as documenting workshops and their open-problem discussions.
\item Supplied literature source, ``Hunt formula for SUq and Uq'': \url{https://eprints.lancs.ac.uk/id/eprint/222139/}. No authors or further bibliographic details are supplied.
\item Supplied literature source, ``Connectedness and Gaussian parts'': \url{https://arxiv.org/abs/2203.08030}. No authors or further bibliographic details are supplied.
\item Literature assessment dated exactly \texttt{2026-08-22}: the supplied curation labels the problem ``partially solved'' and states that L\'evy--Khintchine decompositions are established for \(SU_q\) of arbitrary finite rank and several related quantum groups, while extension to \(q\)-deformations of all simple compact Lie groups is not established. This is supplied dated curation derived from the OWR record, not an independently verified assessment; its time-sensitive status and cited claims require independent checking.
\item The registry labels the statement \texttt{corrected\_verified} and records that the linked report was inspected and its concluding conjecture isolated from a spliced formula and acknowledgements. This is a supplied verification claim, not a claim that this draft's preparer visited or checked the report.
\item The spliced source acknowledgements attribute partial support to an unnamed author from the Polish National Science Center, SONATA grant \texttt{2016/21/D/ST1 /03010}, and the Polish National Agency for Academic Exchange, POLONIUM program \texttt{PPN/BIL/2018/1/00197/U/00021}. These acknowledgements do not identify the conjecture's proposer.
\end{itemize}
\end{document}
